\begin{figure}[htbp]
\centering
\begin{tikzpicture}[
    font=\small,
    actor/.style={draw=#1, fill=#1!7, rounded corners=2pt, align=center,
                  minimum width=2.8cm, minimum height=0.95cm, inner sep=3pt},
    actor/.default=subsectionblue,
    vida/.style={draw=gray!45, dashed},
    ida/.style={-{Stealth[length=2mm]}, thick, draw=gray!70!black},
    vuelta/.style={-{Stealth[length=2mm]}, thick, dashed, draw=gray!70!black},
    msg/.style={font=\scriptsize, above, inner sep=2pt},
    paso/.style={draw=gray!60, fill=white, rounded corners=2pt, font=\scriptsize,
                 align=center, inner sep=3pt},
    nota/.style={font=\scriptsize\itshape, text=ugrred, align=right},
]

% Columnas
\def\N{0}     % navegador
\def\B{4.2}   % backend
\def\L{8.4}   % modelo
\def\H{12.6}  % herramientas y BD

% Tramo posterior a la respuesta
\fill[subsectionblue!6, rounded corners=3pt] (-1.5,-9.5) rectangle (14.1,-13.65);

% Líneas de vida
\foreach \x in {\N,\B,\L,\H} \draw[vida] (\x,-0.5) -- (\x,-13.65);

\node[font=\scriptsize\bfseries, text=subsectionblue, anchor=north east, align=right,
      fill=subsectionblue!6, inner sep=2pt]
     at (14.0,-9.6) {Después de \texttt{fin\_respuesta}:\\el usuario ya está oyendo la respuesta};

% Cabeceras
\node[actor]        at (\N,0) {Navegador};
\node[actor]        at (\B,0) {Backend\\[-1pt]\scriptsize\texttt{agent.py}};
\node[actor=ugrred] at (\L,0) {Modelo\\[-1pt]\scriptsize\texttt{qwen3:8b}};
\node[actor=gray]   at (\H,0) {Herramientas\\[-1pt]\scriptsize y base de datos};

% --- Respuesta ------------------------------------------------------------
\draw[ida]    (\N,-1.0) -- node[msg] {pregunta}                          (\B,-1.0);
\node[paso]   at (\B,-1.8) {atajos deterministas (saldo, Bizum): no aplican};
\draw[ida]    (\B,-2.6) -- node[msg] {prompt + historial + herramientas} (\L,-2.6);
\draw[vuelta] (\L,-3.4) -- node[msg] {pide \texttt{consultar\_movimientos}} (\B,-3.4);
\draw[ida]    (\B,-4.2) -- node[msg] {\texttt{tool\_inicio}, \texttt{sql}} (\N,-4.2);
\draw[ida]    (\B,-5.0) -- node[msg] {SQL en solo lectura}               (\H,-5.0);
\draw[vuelta] (\H,-5.8) -- node[msg] {filas}                             (\B,-5.8);
\draw[ida]    (\B,-6.6) -- node[msg] {resultado de la herramienta}       (\L,-6.6);
\draw[vuelta] (\L,-7.4) -- node[msg] {respuesta en \textit{streaming}}   (\B,-7.4);
\draw[ida]    (\B,-8.2) -- node[msg] {\texttt{texto}, una frase por evento} (\N,-8.2);
\node[nota, anchor=east] at (\N-0.1,-8.2) {la voz\\empieza};
\draw[ida]    (\B,-9.0) -- node[msg] {\texttt{fin\_respuesta}}           (\N,-9.0);

% --- Motor visual ---------------------------------------------------------
\node[paso]   at (\B,-10.1) {¿hay algo que mostrar?\\regla determinista: sí};
\draw[ida]    (\B,-11.0) -- node[msg] {\texttt{visual\_generando}}        (\N,-11.0);
\draw[ida]    (\B,-11.8) -- node[msg] {llamada dedicada (3 filas)} (\L,-11.8);
\draw[vuelta] (\L,-12.5) -- node[msg] {tipo, spec sin datos y motivo}    (\B,-12.5);
\draw[ida]    (\B,-13.25) -- node[msg] {\texttt{grafico} con datos} (\N,-13.25);

\end{tikzpicture}
\caption{Secuencia de un turno con consulta al histórico y gráfico. La voz
empieza con la primera frase del texto. El motor visual solo se ejecuta cuando
la respuesta ya está completa, de modo que la segunda llamada al modelo no
retrasa lo que el usuario oye.}
\label{fig:turno}
\end{figure}
